// Ejercicio 8.1. Longitud manual.
#include <stdio.h>

int main() {
   char palabra[100];
   
   printf("Introduce una palabra: ");
   scanf("%s", palabra);
   
   int i = 0;
   while (palabra[i] != '\0') {
      i++;
   }
   
   printf("La longitud es: %d\n", i);
   return 0;
}   


// Ejercicio 8.2. Contar vocales.
#include <stdio.h>

int es_vocal(char c) {
   c = (c >= 'A' && c <= 'Z') ? c + 32 : c; // convertir a minúscula
   return c == 'a' || c == 'e' || c == 'i' || c == 'o' || c == 'u';
}

int main() {
   char palabra[100];
   int contador = 0;
   
   printf("Palabra: ");
   scanf("%s", palabra);
   
   for (int i = 0; palabra[i] != '\0'; i++) {
      if (es_vocal(palabra[i])) {
         contador++;
      }
   }
   
   printf("Número de vocales: %d\n", contador);
   return 0;
}   


// Ejercicio 8.3. Invertir cadena.
#include <stdio.h>

int main() {
   char original[100], invertida[100];
   
   printf("Introduce una palabra: ");
   scanf("%s", original);
   
   int longitud = 0;
   while (original[longitud] != '\0') {
      longitud++;
   }
   
   for (int i = 0; i < longitud; i++) {
      invertida[i] = original[longitud - 1 - i];
   }
   invertida[longitud] = '\0';
   
   printf("Invertida: %s\n", invertida);
   return 0;
} 


// Ejercicio 8.4. Comparación alfabética.
#include <stdio.h>

int main() {
   char cad1[100], cad2[100];
   printf("Cadena 1: ");
   scanf("%s", cad1);
   printf("Cadena 2: ");
   scanf("%s", cad2);
   
   char ch1 = cad1[0];
   char ch2 = cad2[0];
   
   if(ch1==ch2) {
      printf("Empiezan igual\n");
   } else if(ch1<ch2) {
      printf("cad1 va antes que cad2\n");
   } else {
      printf("cad2 va antes que cad1\n");
   }
   
   return 0;
} 

// Ejercicio 8.5.  Concatenar nombres.
#include <stdio.h>
#include <string.h>

int main() {
   char nombre[50], apellido[50], completo[99];
   
   printf("Introduce nombre: ");
   scanf("%s", nombre);
   
   printf("Introduce apellido: ");
   scanf("%s", apellido);
   
   int ultimo = strlen(nombre);
   for(int i=0; i<ultimo; i++) {
      completo[i] = nombre[i];
   }
   completo[ultimo]=' ';
   for(int i=0; i<strlen(apellido); i++) {
      completo[ultimo+i+1] = apellido[i]; 
   }
   ultimo = strlen(nombre)+strlen(apellido)+1;
   completo[ultimo]='\0';
   
   printf("Nombre completo: %s\n", completo);
   return 0;
}


// Ejercicio 8.6. Contar mayúsculas
#include <stdio.h>

int contar_mayusculas(char cad[]) {
   int contador = 0;
   for(int i=0; cad[i]!='\0'; i++) {
      if(cad[i]>='A' && cad[i]<='Z') {
         contador++;
      }
   }
   return contador;
}
int main() {
   char palabra[40];
   printf("Palabra: ");
   scanf("%s", palabra);
   
   int n = contar_mayusculas(palabra);
   
   printf("%s tiene %d mayúsculas\n", palabra, n);
   return 0;
}


// Ejercicio 8.7. Convertir a minúsculas
#include <stdio.h>

void convertir_minusculas(char s[]) {
   for (int i = 0; s[i] != '\0'; i++) {
      if (s[i] >= 'A' && s[i] <= 'Z') {
         s[i] = s[i] +32;
      }
   }
}

int main() {
   char palabra[40];
   printf("Introduce una palabra: ");
   scanf("%s", palabra);
   
   convertir_minusculas(palabra);
   
   printf("%s\n", palabra);
   return 0;
}   


// Ejercicio 8.8
#include <stdio.h>

int main() {
   char palabra[40];
   char letra;
   printf("Palabra: ");
   scanf("%s", palabra);
   printf("Letra: ");
   scanf(" %c", &letra); // Espacio antes de %c 
   
   int posicion = -1;
   for(int i=0; palabra[i]!='\0'; i++) {
      if(palabra[i]==letra) {
         posicion = i+1;
         break;
      }
   }
   
   if(posicion==-1) {
      printf("La letra %c no está en %s\n", letra, palabra);
   } else {
      printf("Letra %c en posición %d de %s\n", letra, posicion, palabra);
   }
   
   return 0;
}   


\par\vfill\null\pagebreak

// Ejercicio 8.9. Detectar palíndromos.
#include <stdio.h>

int es_palindromo(char* s);

int main() {
   char palabra[100];
   printf("Introduce una palabra: ");
   scanf("%s", palabra);
   
   if (es_palindromo(palabra)) {
      printf("Es un palíndromo.\n");
   } else {
      printf("No es un palíndromo.\n");
   }
   return 0;
}  

int es_palindromo(char* s) {
   int i = 0, j;
   
   while (s[i] != '\0') {
      i++;
   } 
   j = i - 1; i = 0;
   
   while (i < j) {
      if (s[i] != s[j]) {
         return 0;
      } 
      i++; j--;
   }
   return 1;
}


// Ejercicio 8.10. Comparar longitudes
#include <stdio.h>

int main() {
   char palabra_1[40], palabra_2[40];
   printf("Palabra 1: ");
   scanf("%s", palabra_1);
   printf("Palabra 2: ");
   scanf("%s", palabra_2);
   
   int longitud_1=0;
   while(palabra_1[longitud_1]!='\0') {
      longitud_1++;
   }
   int longitud_2=0;


   
   while(palabra_2[longitud_2]!='\0') {
      longitud_2++;
   }   
   if(longitud_1==longitud_2) {
      printf("Miden lo mismo\n");
   } else if(longitud_1>longitud_2) {
      printf("%s es más larga que %s\n", palabra_1, palabra_2);
   } else {
      printf("%s es más corta que %s\n", palabra_1, palabra_2);
   }  
   return 0;
}  


// Ejercicio 8.11. Última aparición de letra
#include <stdio.h>

int main() {
   char palabra[40];
   char letra;
   printf("Palabra: ");
   scanf("%s", palabra);
   printf("Letra: ");
   scanf(" %c", &letra); // Espacio antes de %c
   
   int posicion = -1;
   for(int i=0; palabra[i]!='\0'; i++) {
      if(palabra[i]==letra) {
         posicion = i;
      }
   }
   
   if(posicion == -1) {
      printf("La letra %c no está en %s\n", letra, palabra);
   } else {
      printf("La última %c está en la posición %d\n", letra, posicion);
   }
   return 0;
}  


// Ejercicio 8.12. Reemplazar caracteres
#include <stdio.h>

int main() {
   char palabra[40];
   char letra_1, letra_2;
   printf("Palabra: ");
   scanf("%s", palabra);
   printf("Letra original: "); 
   scanf(" %c", &letra_1); // Espacio antes de %c



   printf("Letra sustituta: ");
   scanf(" %c", &letra_2); // Espacio antes de %c
   
   for(int i=0; palabra[i]!='\0'; i++) {
      if(palabra[i]==letra_1) {
         palabra[i] = letra_2;
      }
   }
 
   printf("%s\n", palabra);  
   return 0;
}  